% !TEX root = ../main.tex
\corrigesChapitre{17}{Ondes mécaniques}
\corrigesSouples

\begin{blocLivre}{cor-17-1}
\begin{corrige}{1}{La ola dans un stade}
\demarche{Une onde est la propagation d'une perturbation dans un milieu, sans transport de matière.}
\begin{questions}
\item Une perturbation (se lever) se \res{\text{propage}} de proche en proche ; le milieu est l'ensemble des \res{\text{spectateurs}}.
\item \res{\text{Non}}, chacun reste à sa place : l'onde transporte une perturbation, de l'\res{\text{énergie}}, pas de matière.
\item \res{\text{Transversale}} : les spectateurs bougent verticalement, la vague avance horizontalement.
\item Longitudinale : le \res{\text{son}} ; transversale : une onde le long d'une \res{\text{corde}}.
\end{questions}
\end{corrige}
\end{blocLivre}
\begin{blocLivre}{cor-17-2}
\begin{corrige}{2}{Le coup de tonnerre}
\demarche{La lumière arrive presque instantanément : le retard mesuré est celui du son, $d=v\,\Delta t$.}
\begin{questions}
\item $d=340\times6{,}0=\res{\qty{2,04e3}{m}}$, soit \qty{2,04}{km}.
\item $\dfrac{2040}{\num{3,0e8}}\approx\res{\qty{6,8e-6}{s}}$ : un millionième de la durée mesurée, \res{\text{négligeable}}.
\item En \qty{3}{s}, le son parcourt $340\times3=\qty{1020}{m}\approx\qty{1}{km}$.
\item Le son, onde mécanique, a besoin d'un milieu matériel : \res{\text{pas dans le vide}}. La lumière, si.
\end{questions}
\end{corrige}
\end{blocLivre}
\begin{blocLivre}{cor-17-3}
\begin{corrige}{3}{Une perturbation le long d'une corde}
\demarche{Repérer le sommet de la perturbation aux deux dates : $v=\dfrac{d}{t_2-t_1}$.}
\begin{questions}
\item $d=3{,}0-1{,}4=\qty{1,6}{m}$ en \qty{0,20}{s} : $v=\dfrac{1{,}6}{0{,}20}=\res{\qty{8,0}{m.s^{-1}}}$.
\item Il reste \qty{1,0}{m} à parcourir : $t=0{,}20+\dfrac{1{,}0}{8{,}0}=\res{\qty{0,325}{s}}$.
\item $\dfrac{0{,}80}{8{,}0}=\res{\qty{0,10}{s}}$.
\item \res{\text{Transversale}} : la corde se déforme perpendiculairement à la propagation.
\end{questions}
\schema[corrige]{17-ex03}
\end{corrige}
\end{blocLivre}
\begin{blocLivre}{cor-17-4}
\begin{corrige}{4}{Lire un oscillogramme}
\demarche{Compter les divisions d'un motif complet, multiplier par la durée d'une division, puis $f=\dfrac{1}{T}$ et $\lambda=\dfrac{v}{f}$.}
\begin{questions}
\item \res{\text{Périodique}} (le motif se répète), mais \res{\text{pas sinusoïdal}}.
\item 2 périodes occupent 8 divisions : $T=4\times\qty{0,5}{ms}=\res{\qty{2,0}{ms}}$ ; $f=\dfrac{1}{\num{2,0e-3}}=\res{\qty{500}{Hz}}$.
\item $\lambda=\dfrac{340}{500}=\res{\qty{0,68}{m}}$.
\item \res{\text{Oui}} : \qty{500}{Hz} est compris entre \qty{20}{Hz} et \qty{20}{kHz}.
\end{questions}
\schema[corrige]{17-ex04}
\end{corrige}
\end{blocLivre}
\begin{blocLivre}{cor-17-5}
\begin{corrige}{5}{Un même son dans trois milieux}
\demarche{La fréquence est fixée par la source ; $\lambda=\dfrac{v}{f}$ change avec le milieu.}
\begin{questions}
\item \res{\text{Non}} : elle est imposée par la source, le diapason.
\item Air : $\dfrac{340}{440}\approx\res{\qty{0,773}{m}}$ ; eau : $\res{\qty{3,41}{m}}$ ; acier : $\res{\qty{13,4}{m}}$.
\item Dans l'\res{\text{acier}} : la perturbation se transmet d'autant plus vite que le milieu est rigide.
\item Le son se propage plus vite dans l'acier que dans l'air : il arrive \res{\text{plus tôt}} par le rail.
\end{questions}
\piege{Croire que la fréquence change avec le milieu : c'est la longueur d'onde qui change, avec la célérité.}
\end{corrige}
\end{blocLivre}
\begin{blocLivre}{cor-17-6}
\begin{corrige}{6}{Simuler une onde périodique, en Python}
\demarche{Le programme calcule $y(x)$ à une date fixée : c'est une « photographie » de l'onde.}
\begin{questions}
\item \code{A} : l'\res{\text{amplitude}} ; \code{T} : la \res{\text{période}} ; \code{lam} : la \res{\text{longueur d'onde}}. Représentation \res{\text{spatiale}} ($t$ fixé, $x$ varie).
\item Les valeurs se répètent tous les \res{\qty{1,0}{m}} : c'est la longueur d'onde.
\item $v=\dfrac{\lambda}{T}=\dfrac{1{,}0}{0{,}50}=\res{\qty{2,0}{m.s^{-1}}}$ et $f=\dfrac{1}{T}=\res{\qty{2,0}{Hz}}$.
\item Fixer \code{x = 0} et faire varier la date dans la boucle, par exemple \code{t = 0.05 * i}.
\end{questions}
\end{corrige}
\end{blocLivre}
\begin{blocLivre}{cor-17-7}
\begin{corrige}{7}{Un séisme}
\demarche{Les deux trains d'ondes parcourent la même distance $d$ avec des célérités différentes : leurs durées de trajet diffèrent de \qty{12}{s}.}
\begin{questions}
\item Ondes P à \res{\qty{16,8}{s}}, ondes S à \res{\qty{28,8}{s}} : les ondes \res{\text{P}} sont les plus rapides (schéma).
\item Retard : $28{,}8-16{,}8=\res{\qty{12}{s}}$.
\item $\dfrac{d}{3{,}5}-\dfrac{d}{6{,}0}=12$, soit $d\times\num{0,119}=12$ : $\res{d\approx\qty{101}{km}}$.
\item Une station donne un cercle de rayon $d$ ; deux cercles se coupent en deux points ; la \res{\text{troisième}} lève l'ambiguïté.
\end{questions}
\schema[corrige]{17-ex07}
\end{corrige}
\end{blocLivre}
\begin{blocLivre}{cor-17-8}
\begin{corrige}{8}{L'écholocation de la chauve-souris}
\demarche{L'écho fait un aller-retour : $d=\dfrac{v\,\Delta t}{2}$ ; puis $\lambda=\dfrac{v}{f}$.}
\begin{questions}
\item \res{\text{Non}} : \qty{40}{kHz} dépasse \qty{20}{kHz}, ce sont des ultrasons.
\item $d=\dfrac{340\times\num{0,012}}{2}=\res{\qty{2,04}{m}}$.
\item $\lambda=\dfrac{340}{\num{40e3}}=\res{\qty{8,5}{mm}}$.
\item Un son grave (\qty{1}{kHz}) a une longueur d'onde de \qty{34}{cm}, bien plus grande qu'un insecte : les ultrasons permettent de détecter de \res{\text{petites proies}}.
\end{questions}
\piege{Oublier le facteur 2 : l'écho parcourt deux fois la distance qui sépare l'émetteur de l'obstacle.}
\end{corrige}
\end{blocLivre}
\begin{blocLivre}{cor-17-9}
\begin{corrige}{9}{Mesurer la célérité des ultrasons}
\demarche{Deux points en phase sont distants d'un nombre entier de longueurs d'onde : dix « retours en phase » valent $10\,\lambda$.}
\begin{questions}
\item Les deux récepteurs vibrent \res{\text{de la même façon au même instant}} : leurs maxima coïncident.
\item $10\,\lambda=\qty{8,5}{cm}$ donc $\lambda=\res{\qty{8,5}{mm}}$.
\item $v=\lambda\,f=\num{8,5e-3}\times\num{40,0e3}=\res{\qty{340}{m.s^{-1}}}$.
\item L'erreur de lecture de la règle se répartit sur dix longueurs d'onde : la \res{\text{précision}} est meilleure.
\end{questions}
\end{corrige}
\end{blocLivre}
\begin{blocLivre}{cor-17-10}
\begin{corrige}{10}{La houle}
\demarche{La distance entre deux crêtes est la longueur d'onde ; $v=\dfrac{\lambda}{T}$.}
\begin{questions}
\item La \res{\text{longueur d'onde}}, $\lambda=\qty{100}{m}$.
\item $v=\dfrac{100}{8{,}0}=\res{\qty{12,5}{m.s^{-1}}}$, soit \res{\qty{45}{km/h}}.
\item Il monte et descend (en décrivant de petites boucles) sans être entraîné : \res{\text{pas de transport de matière}}.
\item $\dfrac{\num{50e3}}{12{,}5}=\qty{4000}{s}\approx\res{\qty{1,1}{h}}$.
\end{questions}
\end{corrige}
\end{blocLivre}
\begin{blocLivre}{cor-17-11}
\begin{corrige}{11}{Une double périodicité}
\demarche{En abscisse : $x$ pour la représentation spatiale, $t$ pour la temporelle. Mesurer un motif complet.}
\begin{questions}
\item À gauche, $y$ en fonction de $x$ : \res{\text{spatiale}} ; à droite, $y$ en fonction de $t$ : \res{\text{temporelle}}.
\item $\res{\lambda=\qty{0,40}{m}}$ et $\res{T=\qty{0,050}{s}}$ (schéma).
\item $v=\dfrac{0{,}40}{\num{0,050}}=\res{\qty{8,0}{m.s^{-1}}}$ ; $f=\dfrac{1}{\num{0,050}}=\res{\qty{20}{Hz}}$.
\item \res{\qty{1}{cm}}, la même sur les deux graphiques : c'est l'amplitude du mouvement des points.
\end{questions}
\schema[corrige]{17-ex11}
\end{corrige}
\end{blocLivre}
\begin{blocLivre}{cor-17-12}
\begin{corrige}{12}{L'échographie}
\demarche{$\lambda=\dfrac{v}{f}$ ; aller-retour : profondeur $=\dfrac{v\,\Delta t}{2}$.}
\begin{questions}
\item $\lambda=\dfrac{1540}{\num{3,5e6}}\approx\res{\qty{0,44}{mm}}$.
\item $p=\dfrac{1540\times\num{65e-6}}{2}\approx\res{\qty{5,0}{cm}}$.
\item Une fréquence plus élevée donne une longueur d'onde plus \res{\text{petite}} : on voit des détails plus fins (mais les ultrasons pénètrent moins).
\item Il chasse l'\res{\text{air}} entre la sonde et la peau, qui réfléchirait presque tous les ultrasons.
\end{questions}
\end{corrige}
\end{blocLivre}
\begin{blocLivre}{cor-17-13}
\begin{corrige}{13}{Le tsunami}
\demarche{$v=\sqrt{g\,h}$ ; la période reste la même, donc $\lambda=v\,T$ suit les variations de la célérité.}
\begin{questions}
\item $v=\sqrt{9{,}81\times4000}\approx\res{\qty{198}{m.s^{-1}}}$, soit \res{\qty{713}{km/h}}, la vitesse d'un avion.
\item $\dfrac{\num{1e6}}{198}\approx\qty{5,05e3}{s}\approx\res{\qty{1,4}{h}}$.
\item $T=\qty{900}{s}$ : au large, $\lambda\approx198\times900\approx\res{\qty{178}{km}}$ ; à la côte, $v\approx\qty{9,9}{m.s^{-1}}$ et $\lambda\approx\res{\qty{8,9}{km}}$.
\item L'onde ralentit et se raccourcit : son énergie se concentre sur une distance vingt fois plus courte, et sa \res{\text{hauteur augmente}} fortement.
\end{questions}
\schema[corrige]{17-ex13-cor}
\end{corrige}
\end{blocLivre}
\begin{blocLivre}{cor-17-14}
\begin{corrige}{14}{La profondeur d'un puits}
\demarche{Sans le son : $h=\dfrac{1}{2}\,g\,t^2$. Avec le son, l'équation $\sqrt{\dfrac{2h}{g}}+\dfrac{h}{v}=t$ se résout par dichotomie.}
\begin{questions}
\item $h=\dfrac{1}{2}\times9{,}81\times2{,}5^2\approx\res{\qty{30,7}{m}}$.
\item La solution est toujours dans l'intervalle, que l'on \res{\text{coupe en deux}} à chaque étape en gardant la moitié qui la contient.
\item $\res{\qty{28,6}{m}}$ : en négligeant le son, on surestime la profondeur d'environ \res{\qty{7}{\percent}}.
\item $t_2=\dfrac{28{,}6}{340}\approx\qty{0,084}{s}$ et $t_1=\sqrt{\dfrac{2\times28{,}6}{9{,}81}}\approx\qty{2,415}{s}$ : la somme vaut bien \qty{2,5}{s}.
\end{questions}
\end{corrige}
\end{blocLivre}
\begin{blocLivre}{cor-17-15}
\begin{corrige}{15}{Le capteur à ultrasons d'un robot}
\demarche{Distance $=\dfrac{v\,\Delta t}{2}$, avec la vraie célérité pour la distance réelle et \qty{340}{m.s^{-1}} pour la distance affichée.}
\begin{questions}
\item $331+0{,}6\times20=\res{\qty{343}{m.s^{-1}}}$ ; à \qty{0}{\degreeCelsius} : \res{\qty{331}{m.s^{-1}}}.
\item Réelle : $\dfrac{331\times\num{5,83e-3}}{2}\approx\res{\qty{0,965}{m}}$ ; affichée : $\dfrac{340\times\num{5,83e-3}}{2}\approx\res{\qty{0,991}{m}}$.
\item $\dfrac{340-331}{331}\approx\res{\qty{2,7}{\percent}}$ : \qty{3}{cm} sur un mètre, sans conséquence pour éviter un obstacle.
\item $\Delta t=\dfrac{2\times2{,}0}{343}\approx\res{\qty{11,7}{ms}}$. Ajouter un capteur de \res{\text{température}} pour corriger la célérité.
\end{questions}
\end{corrige}
\end{blocLivre}
